\documentclass[12pt]{article}
\usepackage[utf8]{inputenc}
\usepackage[T1]{fontenc}
\usepackage{hyperref}
\title{Turoyo Language Resources for K-12 Curriculum Development}
\author{Compiled by Claude Sonnet 5 (claude-sonnet-5) for Andr\'as Kornai}
\date{}
\begin{document}
\maketitle

\section{Introduction}
Turoyo (Surayt; ISO 639-3 tru; Glottocode turo1239) is a Neo-Aramaic language of Tur Abdin in Turkey and of Syria, with a large diaspora. Search found a modern online course series and a text collection, but no comprehensive dictionary.

\section{Primary Educational Resources}
\begin{itemize}
\item \emph{\v{S}lomo Surayt I} (Shabo Talay, 2017), a free online course at the A1--A2 levels developed under the Erasmus+ Aramaic Online Project, and \emph{\v{S}lomo Surayt II} (2020), an advanced course with a glossary of over 2,000 entries: \url{https://bethmardutho.org/?p=1157}
\end{itemize}

\section{Linguistic and Reference Materials}
\begin{itemize}
\item Wikipedia on Turoyo: \url{https://en.wikipedia.org/wiki/Turoyo_language}
\item Omniglot on Turoyo and its scripts: \url{https://www.omniglot.com/writing/turoyo.htm}
\item Glottolog reference entry for Shabo Talay 2017: \url{https://glottolog.org/resource/reference/id/687051}
\item Shabo Talay's publication page: \url{https://fu-berlin.academia.edu/ShaboTalay}
\item FU Berlin Surayt-Aramaic online release notice: \url{https://www.geschkult.fu-berlin.de/en/e/semiarab/semitistik/news/SAOPonlinerelease.html}
\end{itemize}
Talay's 2019 Midyat texts include a glossary limited to words in the texts.

\section{Religious and Community Resources}
A Swedish--Surayt wordlist from a community site: \url{https://www.acsatv.com/?p=668}

\section{Digital and Audio Resources}
The online course is the main digital resource. The site named in search summaries, www.surayt.com, was not itself a search result and is not cited here.

\section{Curriculum Development Context}
Orthographies were agreed at the 2015 Cambridge Surayt conference, and the language can be written in Syriac or Latin script. Course materials exist for adults and older learners. Primary-school materials were not found, and the sources disagree on speaker numbers.
\end{document}